\documentclass[11pt]{article}
\usepackage[margin=1in]{geometry}
\usepackage{amsmath,amssymb,amsthm,mathtools}
\usepackage[hidelinks]{hyperref}
\newtheorem{theorem}{Theorem}
\newtheorem{lemma}[theorem]{Lemma}
\newtheorem{corollary}[theorem]{Corollary}
\title{Effective simplicity for the ordinary exponent-six factorial polynomials}
\author{Internal analytic proof and exact finite certificates; no priority claim}
\date{17 November 2026}
\begin{document}
\maketitle
Throughout, the Gamma shape is one. Put
\[
 P_n(z)=\sum_{k=0}^n\binom nk(6k)!z^{n-k},\qquad
 F_n(t)=\sum_{k=0}^n\binom nk(6k)!t^k.
\]
The overarching goal is to resolve Factorial and/or Strong Factorial.
The result below concerns only the constant-plus-one-variable-sixth-power
two-term family. It does not resolve either unrestricted conjecture.

\begin{theorem}\label{thm:main}
For every integer $n\ge1$, $P_n$ and $F_n$ have only simple complex roots.
The analytic argument proves this for every $n\ge10000$, with every
scaled root in $60<|z/n^6|<240$. The finite interval $1\le n<10000$
is certified exactly by the accompanying modular B\'ezout certificates
and the previously proved real-place parity lemma.
\end{theorem}
\begin{corollary}
The Strong Factorial assertion holds for $A+BY^6$, for all complex $A,B$
with a nonzero polynomial, at every starting index.
\end{corollary}
\begin{proof}[Reduction to the theorem]
$P_n(z)=z^nF_n(1/z)$ has nonzero constant term, and
$P'_{n+1}=(n+1)P_n$. Also
$\mathcal L((A+BY^6)^n)=B^nP_n(A/B)$ when $B\ne0$.
Thus consecutive vanishing moments would give a multiple root of
$P_{n+1}$. The one-term cases are immediate.
\end{proof}

\section{An effective outer enclosure}
Write $a_k=\binom nk(6k)!$, $\theta=k/n$. The entropy bound for a
binomial coefficient and elementary factorial integral bounds give
\[
 \binom nk\le\theta^{-k}(1-\theta)^{-(n-k)},\qquad
 (6k)!\le6ek(6k/e)^{6k},
\]
with the endpoint interpreted by continuity. Hence
\[
 \frac{a_k^{1/k}}{n^6}
 \le(6/e)^6q(\theta)(6ek)^{1/k},\qquad
 q(\theta)=\theta^5(1-\theta)^{-(1-\theta)/\theta}.       \tag{1}
\]
For $0<\theta\le1$,
\[
                          q(\theta)<300/299.           \tag{2}
\]
To see this, the derivative of $\log q$ has numerator
$6\theta+\log(1-\theta)$, which is positive and then negative with
exactly one zero $\theta_*$ in $(0,1)$. It is positive at $299/300$,
since $\log300<299/50$. At its maximum, with $u=1-\theta_*<1/300$,
\[
 \log q(\theta_*)=5\log(1-u)+6u<u<1/300.
\]
Use $e^{1/300}\le300/299$.
For clarity, $\log300<299/50$ follows from the entirely rational bound
$(27/10)^5[1+49/50+(49/50)^2/2]>300$ and $e>27/10$.

For $k\ge1000$, monotonicity of $(6ek)^{1/k}$ gives
$(6ek)^{1/k}<101/100$, because $(101/100)^{1000}>18000>6000e$.
Since $e>19/7$,
\[
 (6/e)^6(300/299)(101/100)
 <(42/19)^6(300/299)(101/100)<120.                      \tag{3}
\]
For $3\le k<1000$ and $n\ge10000$, the cruder coefficient estimate
from the \href{EFFECTIVE_SADDLE_ANNULUS.tex}{saddle-annulus note} gives
\[
 a_k^{1/k}/n^6<(177147/128)(k/n)^5<120.
\]
The cases $k=1,2$ also give $a_k^{1/k}/n^6<120$ for $n\ge10000$;
indeed $720/3^5<120$ and $\sqrt{12!/2}/3^5<120$ already suffice.
Thus $a_k^{1/k}<120n^6$ for every coefficient at these indices. The
geometric-series argument applied to $P_n(z)/z^n$ excludes all
$|z|\ge240n^6$.

\section{An exact decomposition and the large saddle}
Put
\[
 I_n(w)=\int_0^\infty(w+x^6)^ne^{-nx}\,dx
       =K_n(w)+J_n(w),
\]
where $K_n$ integrates over $[0,3]$ and $J_n$ over $[3,\infty)$.
Then $P_n(n^6w)=n^{6n+1}I_n(w)$, so multiplicities are unchanged.
This decomposition is exact, and both pieces are polynomials in $w$.

For $|w|\le250$, there is a unique saddle $s_0=\alpha+i\beta$ with
$|s_0-6|<\delta=1/25$. Indeed
$(6-\delta)^5\delta>250$, so Rouch\'e applies to
$w=s^5(6-s)$. Near this saddle use
\[
 \Phi(s,w)=6\log s+\log(1+w/s^6)-s,\quad
 \Phi_0=\Phi(s_0,w),\quad a=1-5/s_0.
\]
The functions are analytic in $w$ on a neighborhood of the closed disk.
They satisfy
\[
 \Phi_s(s_0,w)=0,\quad \Phi_{ss}(s_0,w)=-a,\quad
 \operatorname{Re}a>1/7,\quad |a|<1/5.                 \tag{4}
\]
Writing $G(x)=6\log x-x$ and $S=\operatorname{Re}\Phi_0$,
\[
 S-G(6)=5\operatorname{Re}\log(1+(s_0-6)/6)
          -\operatorname{Re}(s_0-6)>-1/100.             \tag{5}
\]
This follows from
$|\log(1+v)-v|\le|v|^2/[2(1-|v|)]$ and
$\delta/6+5\delta^2/[72(1-\delta/6)]<1/100$.
The rational lower bound $\log6>179/100$ gives $S>47/10$.

\begin{lemma}\label{lem:J}
For every $n\ge10000$ and $|w|\le250$,
\[
 J_n(w)=U_n(w)(1+\varepsilon_n(w)),\qquad
 U_n=e^{n\Phi_0}\sqrt{2\pi/(na)},\qquad
 |\varepsilon_n|<12/\sqrt n.                           \tag{6}
\]
The relative error is analytic on a neighborhood of that closed disk.
In particular $J_n$ is nonzero there. On $|w|\le240$,
\[
 \left|J_n'/J_n-n\frac{6-s_0}{6w}\right|<1,             \tag{7}
\]
where the phase derivative at $w=0$ is interpreted by its removable limit.
\end{lemma}
\begin{proof}
Deform $[3,\infty)$ to the vertical segment $[3,3+i\beta]$ followed by
$[3+i\beta,\infty+i\beta]$. The integrand is entire, and the right side
of the deformation rectangle vanishes at infinity. Thus all contour
pieces are retained. Set $r=3/4$ and first integrate horizontally over
$s=s_0+t$, $|t|\le r$. Here $|s|>5$ and $|w|\le250$ imply
$|\Phi_{sss}|<1/4$. An explicit series upper bound is
\[
 \frac1{125}\left(12+
 \frac{216x(1+x)}{(1-x)^3}+\frac{108x}{(1-x)^2}
 +\frac{12x}{1-x}\right)<1/4,\quad x=250/5^6.           \tag{8}
\]
Taylor's remainder is therefore at most $|t|^3/24$.
Since $1/14-r/24=9/224$, the error in replacing this local integral
by the local Gaussian is at most
\[
 \frac{e^{nS}}{24n(9/224)^2}<\frac{26e^{nS}}n.
\]
Completing the Gaussian adds at most
$[56/(3n)]e^{nS-9n/224}<19e^{nS-9n/224}/n$.
The complete Gaussian is $U_n$, with $|U_n|>5e^{nS}/\sqrt n$.

It remains to bound the entire complement of the local interval.
For $s=x+i\beta$ and $4\le x\le8$, differentiate the real part of
$\Phi=6\log s+\log(1+w/s^6)-s$. With $q=250/4^6$ and
$|\beta|/x\le1/100$,
\[
 \frac{d^2}{dx^2}\operatorname{Re}\Phi(x+i\beta,w)
 \le-\frac1{x^2}\left[
 \frac{6(1-10^{-4})}{(1+10^{-4})^2}
 -\frac{36q}{(1-q)^2}-\frac{6q}{1-q}\right]
 <-\frac3{x^2}\le-\frac3{64}.                         \tag{9}
\]
This bound uses only $|w/s^6|\le q<1$, so the logarithmic chart is
valid throughout the horizontal segment. Its real first derivative at
$x=\alpha$ vanishes. Consequently, on $[4,8]$,
\[
 \operatorname{Re}\Phi(x+i\beta,w)
 \le S-\frac3{128}(x-\alpha)^2.
\]
Outside $|x-\alpha|\le r$, the loss exceeds $1/100$.
This part contributes at most $4e^{nS-n/100}$.

On $[3,4]$, the upper envelope
$[250+(x+\delta)^6]e^{-x}$ increases: its derivative has the sign of
$u^5(6-u)-250>0$, for $u=x+\delta\in[3+\delta,4+\delta]$.
Its maximum is less than $e^{9/2}$, since
\[
 250+(101/25)^6<(27/10)^8(41/25)<e^{17/2}.
\]
Here $e>27/10$ and $(41/25)^2<27/10$.
This interval contributes at most $e^{nS-n/5}$ by $S>47/10$.
The vertical segment at $3$ contributes at most
$\delta e^{nS-n/2}$, since
$250+(76/25)^6<(27/10)^7<e^7$.

For $x\ge8$, the same chart gives
\[
 \frac d{dx}\operatorname{Re}\Phi(x+i\beta,w)
 \le\frac6{x(1-250/x^6)}-1<-1/5.
\]
By (9), its height at eight is at most
$S-(3/128)(49/25)^2<S-1/20$. Thus this infinite tail contributes
at most $5e^{nS-n/20}/n$.
The absolute error after removing $U_n$ is bounded by $e^{nS}$ times
\[
 \frac{26}{n}+\frac{19}{n}e^{-9n/224}
 +4e^{-n/100}+e^{-n/5}+\frac5n e^{-n/20}
 +\delta e^{-n/2}<\frac{60}{n}\quad(n\ge10000).        \tag{10}
\]
For the last inequality, drop the exponentials on the $1/n$ terms
and use $e^v\ge v^2/2$ on the others; the resulting coefficient is
$50+(80000+50+8/25)/10000<60$. This proves (6).

Analyticity follows since $J_n$ is a polynomial and $U_n$ is analytic
and nonzero on a larger disk. Cauchy's estimate on radius-ten circles
inside $|w|\le250$ gives
$|\varepsilon_n'|\le12/(10\sqrt n)$ on $|w|\le240$.
Furthermore
\[
 \frac{U_n'}{U_n}=n\frac{6-s_0}{6w}-\frac{a'}{2a},\quad
 s_0'=\frac1{6s_0^4(5-s_0)},\quad
 \left|\frac{a'}{2a}\right|<1/1000.
\]
The last bound follows from (4), $|s_0|>6-\delta$ and
$|s_0-5|>1-\delta$. Finally
\[
 \frac1{1000}+
 \frac{12/(10\sqrt n)}{1-12/\sqrt n}<1\quad(n\ge10000),
\]
which proves (7).
\end{proof}

\section{The endpoint and the inner root exclusion}
\begin{lemma}\label{lem:K}
For $n\ge10000$ and $|w|\ge50$,
\[
 K_n(w)=\frac{w^n}{n}(1+D_n(w)),\qquad
 |D_n|\le\frac{720\cdot20^7}{50n^5}+e^{-3n}<1/1000,
 \qquad |D_n'|\le\frac{720\cdot20^7}{2500n^5}<1/1000.   \tag{11}
\]
In particular, $K_n$ has no zero there and
$|K_n'/K_n-n/w|<1$.
\end{lemma}
\begin{proof}
For $0\le x\le3$,
\[
                   \log(1+x^6/50)\le19x/20.            \tag{12}
\]
Indeed $\log(1+x^6/50)/x$ increases: its derivative has the sign of
$h(u)=6u/(1+u)-\log(1+u)$, for $0\le u\le729/50$.
On $[0,5]$, $h'=(5-u)/(1+u)^2\ge0$; above five, use
$6u/(1+u)\ge5$ and $\log(1+u)<3$. Finally
$\log(779/50)<57/20$, as certified by
$(27/10)^2[1+17/20+(17/20)^2/2]>779/50$.

Subtract $n\int_0^\infty e^{-nx}dx=1$ from
$nK_n/w^n$. The inequality
$|(1+v)^n-1|\le n|v|(1+|v|)^{n-1}$ and (12) give
\[
 |D_n|\le\frac{n^2}{50}\int_0^\infty x^6e^{-nx/20}dx+e^{-3n}.
\]
The integral is $720(20/n)^7$. Differentiation under the finite integral
gives the bound for $D_n'$, with $50$ replaced by $2500$.
The logarithmic derivative estimate follows directly.
\end{proof}

For $|w|\le60$, a different absolute bound suffices:
$|K_n(w)|\le3\cdot60^n$. In fact the maximum of
$(60+x^6)e^{-x}$ on $[0,3]$ occurs at an endpoint, because
$x^5(6-x)$ increases there. The endpoint values are $60$ and
$789e^{-3}<60$. On the other hand Lemma~\ref{lem:J} gives
\[
 |J_n(w)|>4e^{nS}/\sqrt n>4\cdot100^n/\sqrt n,
\]
using $S>47/10$ and $e^{47/10}>100$.
Thus $|K_n|<|J_n|$ for $n\ge10000$. For example
$3\sqrt n(3/5)^n<4$ follows at $n=10000$ by an exact rational
comparison and then decreases with the integer $n$.
This excludes every zero with $|w|\le60$.

\section{Every large-index zero is simple}
By the outer enclosure and the inner exclusion, every zero has
$60<|w|<240$. On this entire annulus, both $J_n$ and $K_n$ are nonzero.
At a zero of $I_n=J_n+K_n$, differentiation gives
\[
 I_n'=J_n\left(\frac{J_n'}{J_n}-\frac{K_n'}{K_n}\right).
\]
The leading logarithmic derivative difference is
\[
 n\left(\frac{6-s_0}{6w}-\frac1w\right)
                     =-\frac{ns_0}{6w}.
\]
By Lemmas~\ref{lem:J}--\ref{lem:K},
\[
 \left|\frac{J_n'}{J_n}-\frac{K_n'}{K_n}\right|
 >\frac{5n}{6\cdot240}-2=\frac n{288}-2>0.
\]
Hence $I_n'$ is nonzero at every zero, for every $n\ge10000$.
The exact split at three, rather than an unproved selection of formal
saddle contributions, supplies the global two-contribution decomposition.
All estimates apply on full disks or annuli, not just a bulk root region.

\section{Exact verification of the finite prefix}
The companion checker supplies certificates for $2\le n<10000$;
$n=1$ is linear. There are two types.
For $2\le n<30$, use prime $211>6n$. The full-degree reduction of
$F_n$ has a B\'ezout combination with its derivative of degree at most one.
For $30\le n<10000$, take the largest prime $q\le n$ and $d=n-q$.
The exact sieve verifies $q>6\max(1,d)$ in every case (maximum $d=35$).
The registered shifted-Newton lemma, applied with multiplier one, gives
$\overline F_n=\overline F_d$ and forces every rational square divisor
to reduce to a square divisor of $\overline F_d$, preserving its degree.
The checker supplies a modular B\'ezout combination
\[
                U\overline F_d+V\overline F_d'=G,
                \qquad \deg G\le1.
\]
For clarity, the two local inputs can also be read directly. The shifted
$q$-polygon has a horizontal edge of length $d$ followed by an edge of
length $q$ and reduced denominator $q$. A square divisor cannot use the
latter: its doubled occupancy would have to be at least $2q$. Normalize
a rational square divisor $g$ as a primitive integral polynomial with
constant term one. Its horizontal polygon then makes its leading
coefficient a $q$-unit, so its reduction preserves its degree and its
square divides $\overline F_d$. For the small full-degree reductions,
$211>6n$ makes the leading coefficient of $F_n$ a unit and also preserves
every factor degree.

Every real root of $F_n$ is simple. Indeed
$F_n(t)=\mathbb E(1+tY^6)^n$, for positive exponential $Y$.
For even $n$ this is strictly positive on the real axis; for odd $n$,
$F_n'(t)=n\mathbb E[Y^6(1+tY^6)^{n-1}]>0$ there.
Thus a repeated irreducible rational factor has no real root and has even
degree. Any complex repeated root supplies such a rational factor, since
$F_n$ has rational coefficients and characteristic zero.
A positive square-divisor degree is therefore at least two and contradicts
the B\'ezout identity of degree at most one after degree-preserving
reduction. This establishes the entire prefix, not a sample of indices.

The certificate checker verifies every supplied B\'ezout identity, prime
and degree condition, along with the exact rational constants and symbolic
identities in the analytic argument. The universal large-index evidence
is the contour proof, not a finite scan. Independent replay, formal
verification and external review remain false. No unrestricted SFC,
higher-dimensional Factorial, other exponent, or other Gamma-shape theorem
is claimed by this result.
\end{document}
